\documentclass[11pt,a4paper]{article}

% ------------------------------------------------------------------------------
% Packages
% ------------------------------------------------------------------------------
\usepackage[utf8]{inputenc}
\usepackage[T1]{fontenc}
\usepackage[french]{babel}
\usepackage{lmodern}
\usepackage[a4paper,margin=2.2cm]{geometry}
\usepackage{xcolor}
\usepackage{array}
\usepackage{multirow}
\usepackage{float}
\usepackage{caption}
\usepackage{listings}
\usepackage{tikz}
\usetikzlibrary{babel,arrows.meta}
\usepackage{hyperref}

\hypersetup{colorlinks=true, linkcolor=bleu, urlcolor=bleu}

% ------------------------------------------------------------------------------
% Couleurs
% ------------------------------------------------------------------------------
\definecolor{bleu}{HTML}{1E3A8A}
\definecolor{vert}{HTML}{16A34A}
\definecolor{vertfond}{HTML}{DCFCE7}
\definecolor{rouge}{HTML}{DC2626}
\definecolor{rougefond}{HTML}{FEE2E2}
\definecolor{orange}{HTML}{EA580C}
\definecolor{orangefond}{HTML}{FFEDD5}
\definecolor{gris}{HTML}{64748B}
\definecolor{codefond}{HTML}{F8FAFC}
\definecolor{codecadre}{HTML}{CBD5E1}

% ------------------------------------------------------------------------------
% Style du code C++
% ------------------------------------------------------------------------------
\lstset{
  language=C++,
  basicstyle=\ttfamily\footnotesize,
  keywordstyle=\color{bleu}\bfseries,
  commentstyle=\color{gris}\itshape,
  stringstyle=\color{vert},
  backgroundcolor=\color{codefond},
  frame=single,
  rulecolor=\color{codecadre},
  breaklines=true,
  columns=fullflexible,
  keepspaces=true,
  showstringspaces=false,
  numbers=left,
  numberstyle=\tiny\color{gris},
  captionpos=b,
}
\renewcommand{\lstlistingname}{Code}

% ------------------------------------------------------------------------------
% Styles TikZ
% ------------------------------------------------------------------------------
\tikzset{
  case/.style={rectangle, draw=black!60, thick, minimum width=17mm, minimum height=7mm,
               inner sep=0pt, font=\small\bfseries, fill=white},
  garde/.style={case, draw=orange, fill=orangefond},
  neuf/.style={case, draw=vert, fill=vertfond},
  noeud/.style={rectangle, rounded corners=2pt, draw=black!60, thick, minimum width=14mm,
                minimum height=7mm, inner sep=0pt, font=\small\bfseries, fill=white},
  noeudgarde/.style={noeud, draw=orange, fill=orangefond},
  noeudneuf/.style={noeud, draw=vert, fill=vertfond},
  lien/.style={draw=black!45, thick},
  fleche/.style={-{Stealth[length=2.5mm]}, thick, orange},
  indice/.style={font=\scriptsize, text=gris},
}

\setlength{\parindent}{0pt}
\setlength{\parskip}{5pt}
\renewcommand{\arraystretch}{1.2}
\newcommand{\entete}[1]{\textbf{\begin{tabular}[c]{@{}c@{}}#1\end{tabular}}}
\newcommand{\oui}{\textcolor{vert}{\textbf{valide}}}
\newcommand{\non}{\textcolor{rouge}{\textbf{invalide}}}

\title{\textbf{Insertion et validité des itérateurs}\\[4pt]
       \large \texttt{std::vector}, \texttt{std::map} et la table de hachage
       \texttt{std::unordered\_map}, avec des paires (clé, valeur)}
\author{Zakariae Bakkari\\
        \small Programmation Objet Avancée et Design Patterns\\
        \small\url{https://github.com/zakariae-bakkari/cpp_ilisi2}}
\date{13 novembre 2026}

\begin{document}

\maketitle

% ==============================================================================
\section{Les questions}
% ==============================================================================

\begin{enumerate}
  \item Une insertion pendant qu'on garde un itérateur va-t-elle le perturber ?
  \item L'insertion dans un \texttt{vector} et dans une \texttt{map} change-t-elle
        l'itérateur différemment ? Faut-il en reprendre un autre ?
  \item Concrètement : on a un \texttt{vector}, on insère 0, l'itérateur reste-t-il
        valide ? Et pour la \texttt{map} ? Pourquoi ?
  \item Que se passe-t-il avec une table de hachage (\texttt{unordered\_map}) ?
  \item Avec des paires (clé, valeur), dans trois scénarios : données
        \textbf{aléatoires}, \textbf{triées} et \textbf{presque triées}.
\end{enumerate}

\textbf{Code source} (option \textbf{17} du menu \texttt{./run.sh}, ou
\texttt{make run\_16}) :\\
{\scriptsize\url{https://github.com/zakariae-bakkari/cpp_ilisi2/blob/main/iterators/src/16_invalidation_iterateurs_map.cpp}}

% ==============================================================================
\section{La réponse en un tableau}
% ==============================================================================

Un itérateur est \textbf{invalide} quand il ne désigne plus l'élément qu'il désignait. Le
lire ou l'avancer est alors un \textbf{comportement indéfini} : le programme peut afficher
une mauvaise valeur, planter, ou sembler marcher par hasard.

\begin{table}[H]
\centering
\small
\begin{tabular}{|p{3.6cm}|p{4.2cm}|p{3.4cm}|p{3.4cm}|}
\hline
\textbf{Opération} & \textbf{\texttt{vector}} & \textbf{\texttt{map}} & \textbf{\texttt{unordered\_map}}\newline\textbf{(table de hachage)} \\
\hline
insertion, sans réallocation ni rehash
  & \non{} à partir du point d'insertion (\oui{} avant)
  & \oui{}
  & \oui{} \\
\hline
insertion, avec réallocation (\texttt{vector}) ou rehash (\texttt{unordered\_map})
  & \non{} partout
  & (n'arrive jamais)
  & itérateurs \non{},\newline références \oui{} \\
\hline
suppression d'un élément
  & \non{} à partir du point de suppression
  & \oui{} sauf l'élément supprimé
  & \oui{} sauf l'élément supprimé \\
\hline
\end{tabular}
\caption{Ce que deviennent les itérateurs gardés avant l'opération (règles de la norme C++).}
\end{table}

\textbf{Donc oui, il faut parfois reprendre un autre itérateur} : toujours avec un
\texttt{vector} après une insertion au milieu ou une réallocation, et avec une
\texttt{unordered\_map} après un rehash. Avec une \texttt{map}, jamais.

% ==============================================================================
\section{Pourquoi ? La façon dont chaque conteneur range ses éléments}
% ==============================================================================

\subsection{\texttt{vector} : un seul tableau, les éléments collés}

Un \texttt{vector} range ses éléments \textbf{côte à côte} dans un seul bloc de mémoire. Un
itérateur est en fait une adresse dans ce bloc. Insérer 0 au début peut casser cette
adresse de deux façons :

\begin{figure}[H]
\centering
\begin{tikzpicture}[x=17mm, y=10mm]
  % avant
  \node[font=\small\bfseries, anchor=east] at (-0.6,0) {avant};
  \node[case]  at (0,0) {(10,100)};
  \node[garde] (g) at (1,0) {(20,200)};
  \node[case]  at (2,0) {(30,300)};
  \node[case, draw=black!25, text=black!30] at (3,0) {libre};
  \draw[fleche] (1,1) node[above, font=\small, text=orange] {it} -- (g.north);
  % decalage
  \node[font=\small\bfseries, anchor=east] at (-0.6,-2) {après};
  \node[neuf]  at (0,-2) {(0,0)};
  \node[garde] (h) at (1,-2) {(10,100)};
  \node[case]  at (2,-2) {(20,200)};
  \node[case]  at (3,-2) {(30,300)};
  \draw[fleche] (1,-1) node[above, font=\small, text=orange] {it} -- (h.north);
  \node[font=\small, text width=50mm, anchor=west] at (3.7,-2)
       {\textbf{(a) il reste de la place :} tout est décalé d'une case ; \texttt{it} désigne
        maintenant (10,100).};
  % reallocation
  \node[font=\small\bfseries, anchor=east] at (-0.6,-4) {après};
  \node[neuf]  at (0,-4) {(0,0)};
  \node[case]  at (1,-4) {(10,100)};
  \node[case]  at (2,-4) {(20,200)};
  \node[case]  at (3,-4) {(30,300)};
  \node[font=\small, text width=50mm, anchor=west] at (3.7,-4)
       {\textbf{(b) tableau plein (capacité 3) :} un nouveau tableau, plus grand, ailleurs ;
        l'ancien est libéré et \texttt{it} pointe dans le vide.};
\end{tikzpicture}
\caption{Insérer 0 au début d'un \texttt{vector} : dans les deux cas, \texttt{it} est invalide.}
\end{figure}

\textbf{Réponse à « on insère 0, l'itérateur reste-t-il valide ? » : non}, jamais quand on
insère avant lui. Seul un \texttt{push\_back} sans réallocation laisse valides les
itérateurs déjà pris (sauf \texttt{end()}).

\subsection{\texttt{map} : un nœud par paire, reliés par des pointeurs}

Une \texttt{map} est un arbre rouge-noir : chaque paire (clé, valeur) vit dans \textbf{son
propre nœud}, alloué une fois pour toutes. Insérer une paire crée un nouveau nœud et
change quelques pointeurs (et des couleurs, et parfois une rotation), mais \textbf{aucun
nœud existant ne change d'adresse}. Un itérateur de \texttt{map} pointe vers un nœud : il
reste donc valide.

\begin{figure}[H]
\centering
\begin{tikzpicture}[level distance=11mm,
  level 1/.style={sibling distance=34mm}, level 2/.style={sibling distance=17mm},
  edge from parent/.style={lien}]
  \node[noeud] {(20,200)}
    child { node[noeud] {(10,100)}
      child { node[noeudneuf] {(0,0)} }
      child[missing] }
    child { node[noeud] {(30,300)} };
  \node[font=\small, text width=62mm, anchor=west] at (2.8,-1.1)
       {Insérer (0,0) ajoute un nœud (en vert) sous (10,100). Les nœuds (10,100), (20,200)
        et (30,300) restent à la même adresse : un itérateur sur (20,200) reste \oui.};
\end{tikzpicture}
\caption{Insérer 0 dans une \texttt{map} : seuls des pointeurs changent.}
\end{figure}

\subsection{Table de hachage (\texttt{unordered\_map}) : des cases, et des nœuds rangés dedans}

Une table de hachage a un tableau de \textbf{cases} (\textit{buckets}). La clé est hachée
pour choisir sa case, et chaque paire est, comme dans la \texttt{map}, dans son propre
nœud. Quand il y a trop de paires par case, la table fait un \textbf{rehash} : elle prend
plus de cases et y redistribue les nœuds.

\begin{itemize}
  \item Les \textbf{itérateurs} deviennent \non{} après un rehash : l'ordre de parcours des
        cases a changé.
  \item Les \textbf{références et pointeurs} vers une paire restent \oui{} : le nœud
        lui-même n'a pas bougé, il a seulement changé de case.
  \item Sans rehash (par exemple après \texttt{reserve(n)}), une insertion n'invalide aucun
        itérateur.
\end{itemize}

% ==============================================================================
\section{Démonstration (programme 16, partie 1)}
% ==============================================================================

Le programme garde un itérateur sur la paire (20,\,200), fait l'opération, puis dit si
l'itérateur est encore valide. Il ne lit \textbf{jamais} un itérateur invalide (ce serait un
comportement indéfini) : il compare les adresses avant et après, et relit les cases avec
\texttt{v[i]}.

\begin{table}[H]
\centering
\small
\begin{tabular}{|c|p{5.4cm}|p{5.2cm}|p{2.6cm}|}
\hline
& \textbf{Opération} & \textbf{Ce qu'on observe} & \textbf{Itérateur} \\
\hline
1 & \texttt{vector} plein (capacité 3), insertion de (0,0) au début
  & capacité 3 $\to$ 6, tableau déplacé ailleurs & \non \\
\hline
2 & \texttt{vector} avec de la place (capacité 10), insertion de (0,0) au début
  & même tableau, mais la case de \texttt{it} contient (10,100) & \non \\
\hline
3 & \texttt{vector} avec de la place, \texttt{push\_back} de (40,400)
  & rien ne bouge avant la fin, \texttt{*it} = (20,200) & \oui \\
\hline
4 & \texttt{map}, insertion de (0,0) puis de 100\,000 paires
  & \texttt{it->first} = 20, \texttt{it->second} = 200 & \oui \\
\hline
5 & \texttt{unordered\_map} (table de hachage), insertion de (0,0) puis de 100\,000 paires
  & 13 $\to$ 172\,933 cases (rehash) ; la paire est restée à la même adresse
  & \non{} (référence \oui) \\
\hline
6 & \texttt{map}, \texttt{erase(10)}
  & \texttt{it->first} = 20 & \oui \\
\hline
\end{tabular}
\caption{Résultats de la partie 1 du programme 16.}
\end{table}

% ==============================================================================
\section{Les trois scénarios avec 100\,000 paires (programme 16, partie 2)}
% ==============================================================================

On insère 100\,000 paires (clé, valeur $= 10 \times$ clé), les clés de 1 à 100\,000 étant
\textbf{aléatoires}, \textbf{triées}, ou \textbf{presque triées} (triées puis 10\,\% des
valeurs échangées au hasard). On garde un itérateur sur la première paire insérée, et on
regarde à la fin s'il est encore valide. Quatre façons de ranger les paires :

\begin{itemize}
  \item \textbf{\texttt{vector} trié} : chaque paire est insérée à sa place
        (\texttt{lower\_bound} puis \texttt{insert}), pour pouvoir chercher par dichotomie ;
  \item \textbf{\texttt{vector} (\texttt{push\_back})} : les paires sont ajoutées à la fin,
        sans ordre (pas de recherche rapide possible) ;
  \item \textbf{\texttt{map}} : l'arbre rouge-noir de la STL ;
  \item \textbf{\texttt{unordered\_map}} : la \textbf{table de hachage} de la STL.
\end{itemize}

\begin{table}[H]
\centering
\footnotesize
\setlength{\tabcolsep}{3pt}
\begin{tabular}{|l|l|r|r|r|r|c|c|}
\hline
\textbf{Données} & \textbf{Structure} & \entete{Insertion} & \entete{Recherche\\des $n$ clés}
  & \entete{Invali-\\dations} & \entete{Éléments\\décalés}
  & \entete{Itérateur\\gardé} & \entete{Référence\\gardée} \\
\hline
\multirow{4}{*}{aléatoires}
  & \texttt{vector} trié        & 616,7 ms & 8,6 ms  & 18 & 2\,502\,036\,049 & \non & \non \\
  & \texttt{vector} (push\_back) & 0,1 ms   & --      & 18 & 0               & \non & \non \\
  & \texttt{map} (arbre)     & 14,0 ms  & 14,6 ms & 0  & 0               & \oui & \oui \\
  & \textbf{table de hachage} & 4,7 ms   & 0,3 ms  & 13 & 0               & \non & \oui \\
\hline
\multirow{4}{*}{triées}
  & \texttt{vector} trié        & 1,9 ms   & 4,0 ms  & 18 & 0               & \non & \non \\
  & \texttt{vector} (push\_back) & 0,1 ms   & --      & 18 & 0               & \non & \non \\
  & \texttt{map} (arbre)     & 7,6 ms   & 5,5 ms  & 0  & 0               & \oui & \oui \\
  & \textbf{table de hachage} & 2,5 ms   & 0,3 ms  & 13 & 0               & \non & \oui \\
\hline
\multirow{4}{*}{presque triées}
  & \texttt{vector} trié        & 132,3 ms & 5,0 ms  & 18 & 576\,523\,734   & \non & \non \\
  & \texttt{vector} (push\_back) & 0,2 ms   & --      & 18 & 0               & \non & \non \\
  & \texttt{map} (arbre)     & 10,3 ms  & 8,0 ms  & 0  & 0               & \oui & \oui \\
  & \textbf{table de hachage} & 3,2 ms   & 0,3 ms  & 13 & 0               & \non & \oui \\
\hline
\end{tabular}
\caption{100\,000 paires. Temps : médianes sur 7 exécutions (Intel Core i5-1135G7,
\texttt{g++} 15.2, \texttt{-O3}) ; les compteurs sont identiques à chaque exécution.
\og Invalidations \fg{} : nombre de réallocations (\texttt{vector}) ou de rehash
(table de hachage, \texttt{unordered\_map}), chacune invalide tous les itérateurs. \og Éléments décalés \fg{} : total des paires déplacées d'une case par les insertions au
milieu du \texttt{vector}.}
\end{table}

\textbf{Ce qu'on en retient.}
\begin{enumerate}
  \item \textbf{La validité ne dépend pas de l'ordre des données.} Dans les trois scénarios,
        l'itérateur gardé est invalide pour le \texttt{vector} (18 réallocations : la
        capacité double à chaque fois, $2^{17} \geq 100\,000$) et pour
        la table de hachage (13 rehash), et toujours valide pour la \texttt{map}.
  \item \textbf{L'ordre des données change surtout le coût du \texttt{vector} trié.} Avec
        des clés aléatoires, chaque insertion décale en moyenne la moitié du tableau :
        2,5 milliards de déplacements ($\approx n^2/4$), 617~ms. Avec des clés triées, on
        insère toujours à la fin : 0 décalage, 1,9~ms. Presque triées : entre les deux.
  \item \textbf{La \texttt{map} est la seule à garder tous ses itérateurs}, quel que soit
        l'ordre des données, pour un coût d'insertion stable (8 à 14~ms).
  \item \textbf{La table de hachage est la plus rapide pour chercher} (0,3~ms pour
        100\,000 recherches, contre 5 à 15~ms pour la \texttt{map}), mais ses itérateurs ne
        survivent pas aux rehash et ses clés ne sont pas triées.
  \item Les temps de ce rapport et ceux du rapport sur l'arbre rouge-noir n'ont pas été
        mesurés au même moment (la machine allait environ 3 fois plus vite ici) : on ne
        compare que les valeurs d'un même tableau.
\end{enumerate}

% ==============================================================================
\section{Comment éviter le problème}
% ==============================================================================

\begin{itemize}
  \item \textbf{Reprendre l'itérateur renvoyé par l'opération} : \texttt{insert} et
        \texttt{erase} renvoient un itérateur valide.
\begin{lstlisting}
it = v.insert(it, {0, 0});   // it designe maintenant (0,0), et il est valide
it = v.erase(it);            // it designe l'element qui suit l'element supprime
\end{lstlisting}
  \item \textbf{\texttt{vector}} : réserver la place à l'avance (\texttt{v.reserve(n)})
        évite les réallocations, ou garder un \textbf{indice} plutôt qu'un itérateur.
  \item \textbf{\texttt{unordered\_map}} : \texttt{u.reserve(n)} prépare assez de cases pour
        $n$ paires et évite les rehash ; sinon, garder une \textbf{référence} vers la paire
        (elle reste valide) ou simplement la \textbf{clé}.
  \item \textbf{\texttt{map}} : rien à faire, ses itérateurs restent valides (sauf celui de
        l'élément supprimé).
\end{itemize}

% ==============================================================================
\section{Conclusion}
% ==============================================================================

Oui, une insertion peut perturber un itérateur, et tout dépend de la façon dont le
conteneur range ses éléments. Le \texttt{vector} colle ses éléments dans un seul tableau :
insérer 0 au début les décale ou les déplace tous, donc l'itérateur est invalide et il faut
en reprendre un autre. La \texttt{map} met chaque paire dans un nœud qui ne bouge jamais :
ses itérateurs restent valides. La table de hachage est entre les deux : ses nœuds ne
bougent pas (les références restent valides), mais un rehash invalide ses itérateurs. Pour
des paires (clé, valeur), la \texttt{map} est donc le choix sûr quand on garde des
itérateurs, et l'\texttt{unordered\_map} le choix rapide quand on cherche par clé sans
besoin d'ordre.

\end{document}
